\SourceLeaf{269}%
the distance from the middle parallel and also contains the flattening, or the square of \(e\), as a factor.

This advantage can be increased considerably by choosing a slightly different form of the arbitrary function that is almost as convenient for computation. Introducing a second constant \(\alpha\), we set
\[
f\upsilon=\alpha\upsilon-i\log k
\]
By choosing the two constants suitably, we can then make the deviation of the scale factor \(m\) from 1 a quantity of the third order in the distance from the middle parallel, apart from the factor \(ee\), which remains present.

\Article{3}
\indent With this choice of the function \(f\), formula 1 gives the following.
\[
T=\alpha t \quad .\ .\ .\ .\ .\ .\ .\ .\ .\ .\ .\ .\ .\ .\ .\ .\ .\quad(2)
\]
\[
\operatorname{tang}\tfrac{1}{2}U=k\operatorname{tang}\tfrac{1}{2}w^\alpha(\tfrac{1+e\cos w}{1-e\cos w})^{\frac{1}{2}\alpha e} \quad .\ .\ .\ .\ .\ .\ .\ .\ .\quad(3)
\]
For \(m\), the principles developed in the paper already mentioned several times readily give the expression
\[
m=\frac{\alpha A\sin U\BareRoot{(1-ee\cos w^2)}}{a\sin w} \quad .\ .\ .\ .\ .\ .\ .\ .\ .\ .\quad(4)
\]
where \(a\) denotes the semimajor axis of the ellipsoid and \(A\) the radius of the sphere.

Differentiating equation 3 written in logarithmic form gives
\[
\frac{\mathrm{d}U}{\sin U}=\frac{\alpha\mathrm{d}w}{\sin w}-\frac{\alpha ee\sin w\mathrm{d}w}{1-ee\cos w^2}
\]
or
\[
\frac{\mathrm{d}U}{\mathrm{d}w}=\frac{\alpha(1-ee)\sin U}{(1-ee\cos w^2)\sin w} \quad .\ .\ .\ .\ .\ .\ .\ .\ .\ .\ .\ .\quad(5)
\]
Likewise, differentiating equation 4 gives
\[
\begin{aligned}
\mathrm{d}\log m&=\operatorname{cotg}U\mathrm{d}U-\operatorname{cotg}w\mathrm{d}w+\frac{ee\cos w\mathord{.}\sin w\mathrm{d}w}{1-ee\cos w^2}\\
&=\operatorname{cotg}U\mathrm{d}U-\frac{(1-ee)\cos w\mathrm{d}w}{(1-ee\cos w^2)\sin w}
\end{aligned}
\]
